\section{Introduction}

Destructive interference can make a collective spin excitation invisible
to a scattering probe. In cluster magnets, amplitudes from sites within
a unit cell combine into a momentum-dependent form factor.
Its zeros suppress intensity even when the corresponding excitation
remains in the spectrum~\cite{waldmann2003interference,haraldsen2011pentamer,tabrizi2021selection,tabrizi2021anisotropic}.
Weakly coupled clusters provide a useful setting for studying this
separation between excitation energies and their visibility.
Their local ground multiplets form effective spin chains, while
intracell interference organizes the folded spectral
weight~\cite{cheng2022trimer,li2026tetramer}.

Near a form-factor zero, corrections to the scattering operator become
essential. Such corrections modify dimer structure
factors~\cite{mueller2003structure} and introduce tensor contributions
when cluster multiplets mix~\cite{waldmann2005manyspin}.
Symmetry constrains these amplitudes but can leave their relative
strengths dependent on the Hamiltonian.
For isotropic clusters, repeated occurrences of a transition-operator
irrep permit independent reduced matrix elements and their
interference~\cite{tabrizi2021selection}.
This distinction also underlies the separation between systematic
extinctions and parameter-dependent zeros~\cite{tabrizi2026extinctions}.
An exchange parameter that cancels one amplitude therefore need not
cancel every transition involving the same symmetry labels.

Thermal occupation adds a further constraint on cancellation.
Excited initial states contribute additional
transitions, as described by effective-observable corrections and
thermal spectral-weight transfer in spin
ladders~\cite{exius2010vertex,honecker2016thermal}.
Thermal lifting of accidental zeros also occurs in magnetovibrational
interference~\cite{raymond2018interference}.
Here we ask whether exchange retuning can compensate the additional
amplitudes and keep the entire low-energy thermal response dark.
The answer requires both the corrected probe and the space of
amplitudes that physical controls can adjust.

We address this question in a Heisenberg pentamer chain whose local
ground multiplet can be a doublet or a quartet.
The doublet permits cancellation of the entire first-order probe on a
uniform reflection-symmetric periodic chain.
At a specified quartet root, a composite vector instead produces
bond-spin amplitudes that remain incompatible with simultaneous
thermal cancellation.
We establish this obstruction for every finite even ring with at least
four cells, including independent first-order SU(2)-invariant controls
inside each cell.
A four-cell ring makes the physical consequence explicit: its leading
ground-state response can vanish, while thermal occupation restores
two calculable groups of spectral lines.
Microscopic spectra test these predictions, and momentum-window
averaging identifies the background that accompanies finite acceptance.

\section{Model and low-energy observable}

Each cell contains five spin-$1/2$ sites at unit spacing. With $J_1=1$,
$\alpha=J_2/J_1$, and $\beta=J_3/J_1>0$, the Hamiltonian is
\begin{align}
 H_0&=\sum_{n=0}^{N-1}\big[
 \alpha(\bm S_{n,1}\!\cdot\!\bm S_{n,2}
       +\bm S_{n,4}\!\cdot\!\bm S_{n,5})\notag\\
 &\hspace{29mm}+\bm S_{n,2}\!\cdot\!\bm S_{n,3}
       +\bm S_{n,3}\!\cdot\!\bm S_{n,4}\big],\\
 H&=H_0+\beta V,\qquad
 V=\sum_{\langle n,n+1\rangle}\bm S_{n,5}\!\cdot\!\bm S_{n+1,1}.
 \label{eq:rank-ham}
\end{align}
The bond sum is specified as open or periodic in each calculation, and the
microscopic length is $L=5N$. We assume that the isolated cell has a single
ground multiplet of spin $s$, separated from all other multiplets by
$\Dloc>0$. The numerical examples below have $s=1/2$ or $3/2$.

Rotational symmetry within this irrep fixes the local projection,
\begin{equation}
 P_n\bm S_{n,a}P_n=c_a\bm T_n,\qquad
 c_a=c_{6-a},\qquad \sum_a c_a=1.
 \label{eq:rank-projection}
\end{equation}
Here $\bm T_n$ is a spin-$s$ operator. The leading Hamiltonian is an
antiferromagnetic effective chain with exchange $\beta c_1^2$.
The scattering intensity additionally depends on how the microscopic
probe projects into this effective spin space.

It is convenient to center the intra-cell Fourier phase:
\begin{align}
 B_n^z(q)&=\sum_{a=1}^5 e^{iq(a-3)}S_{n,a}^z,\\
 O_q&=\frac{1}{\sqrt{5N}}\sum_n e^{iQn}B_n^z(q),
 \quad Q=5q\pmod{2\pi}.
\end{align}
Changing the coordinate origin adds only a global phase to $O_q$ and does not
change a spectral weight. Equation~\eqref{eq:rank-projection} gives
\begin{align}
 P O_qP&=A(q,\alpha)D_Q,\notag\\
 A(q,\alpha)&=\sum_a c_a\cos[q(a-3)],\\
 D_Q&=\frac{1}{\sqrt{5N}}\sum_n e^{iQn}T_n^z.
 \label{eq:rank-amplitude}
\end{align}
The normalization of $D_Q$ is fixed by the microscopic probe and is
kept identical across all comparisons.

At zero temperature, we use the integrated inelastic low-band weight
\begin{equation}
 \Wlow=\frac{1}{n_G}\sum_{r=1}^{n_G}
 \sum_{0<E_m-E_G<\Dloc/2}|\langle m|O_q|G_r\rangle|^2.
 \label{eq:rank-weight}
\end{equation}
All transitions within the ground eigenspace are excluded.
The gap $\Dloc$ in this cutoff and in $g=\beta/\Dloc$ is evaluated at the
unshifted local root. For each finite system, the cutoff contains the
complete retained low band in the weak-coupling regime used below.
The microscopic expansion is taken at fixed $N$; band isolation is
checked separately for each numerical calculation.
